% ================================================================
% ==== FILE: content/cycles/cycles_k5.tex
% ================================================================

\subsubsection{Zyklen auf \texorpdfstring{$K_5$}{K_5}}
\label{sec:cycles-k5}

Ebene$K_5$entspricht dem frühen neuronalen (bioelektrischen) Kontinuum.
Es wird durch das Vorhandensein einer erregbaren Membran, eines Ionenkanals, definiert
Architektur, schwellenwertgesteuerte Spitzen, Erholungsdynamik und
elektrisch-chemische Kopplung.
Zyklen weiter$K_5$stellen die grundlegenden wiederkehrenden Prozesse dar, die beschenken
das System mit Erregbarkeit, Signalisierung, Rhythmik und Proto-Netzwerk
Koordinierung.

\subsubsection{Überblick über die Zyklusstruktur auf \texorpdfstring{$K_5$}{K_5}}

According to the representability theorem for neurons (the corresponding derivation note),
the state space $\Omega(K_5)$ includes:
\[
  V_m(t),\;
  P_{\mathrm{ion}}(t),\;
  \{g_i(t)\},\;
  \{w_{ij}(t)\},\;
  \mathrm{Ca}^{2+}(t),\;
  \mathrm{channel\ states},
\]
mit Schwellenwerten:
\[
  \Theta_{\mathrm{spike}},\;
  \Theta_{\mathrm{exc}},\;
  \Theta_{\mathrm{noise-elec}},\;
  \Theta_{\mathrm{front}},\;
  \Theta_{\mathrm{refrac}},\;
  \Theta_{\mathrm{spatial}}.
\]

Zyklen entstehen, wenn Ströme entstehen$J(t)$, Potenziale$P(t)$, und Achsen$A(t)$
return to an equivalent configuration after a finite time $\tau$.
Die Zeit existiert weiter$K_5$nur weil diese Zyklen stabil und wiederkehrend sind.

Die Kernfamilien der Zyklen sind:

\begin{itemize}
  \item Aktionspotentialzyklen$C_{\mathrm{spike}}$,
  \item Oszillationszyklen unterhalb der Schwelle$C_{\mathrm{sub}}$,
  \item Refraktär-Reset-Zyklen$C_{\mathrm{refrac}}$,
  \item Channel-Gating-Zyklen$C_{\mathrm{gate}}$,
  \item Kalzium/Second-Messenger-Zyklen$C_{\mathrm{Ca}}$,
  \item synaptische Potenzierungs-/Depressionszyklen$C_{\mathrm{syn}}$,
  \item rhythmische Zyklen vernetzen$C_{\mathrm{net}}$das erweiterte Kontinuum bilden$K_5'$.
\end{itemize}

\subsubsection{1. Aktionspotenzialzyklen$C_{\mathrm{spike}}$}

Der Spike ist ein Phasenübergang, der verursacht wird durch:
\[
  T(t) > \Theta_{\mathrm{spike}},
\]
as stated in the corresponding derivation note.

Ein vollständiger Zyklus besteht aus:
\[
  \mathrm{rest} 
  \;\to\;
  \mathrm{depolarisation}
  \;\to\;
  \mathrm{overshoot}
  \;\to\;
  \mathrm{repolarisation}
  \;\to\;
  \mathrm{hyperpolarisation}
  \;\to\;
  \mathrm{rest}.
\]

Formal:
\[
  C_{\mathrm{spike}}
  =\{(V_m(t), g_i(t), P_{\mathrm{ion}}(t))\}_{t_0}^{t_0+\tau_{\mathrm{spike}}},
\]
with period $\tau_{\mathrm{spike}} < \infty$.

Der Spike-Zyklus ist der kanonische elektrische Zyklus$K_5$und die
zentrales Element seiner strukturellen Identität.

\subsubsection{2. Unterschwellige Oszillationszyklen$C_{\mathrm{sub}}$}

Wenn das Membranpotential oszilliert, ohne sich zu kreuzen
$\Theta_{\mathrm{spike}}$ but approaches the excitable regime:
\[
  |V_m(t) - V_{\mathrm{rest}}| < \Theta_{\mathrm{exc}},
\]
es entstehen immer wiederkehrende Schwingungen:
\[
  V_m(t) = V_{\mathrm{rest}} + A\sin(\omega t).
\]

Diese Zyklen sind Vorläufer der Rhythmogenese$K_5'$.

\subsubsection{3. Refraktär-Reset-Zyklen$C_{\mathrm{refrac}}$}

Following a spike, thresholds shift (the corresponding derivation note):
\[
  \Theta_{\mathrm{spike}}(t+\delta t)
  >
  \Theta_{\mathrm{spike}}(t).
\]

Der Refraktärzyklus ist:
\[
  \mathrm{spike} 
  \to 
  \mathrm{absolute\ refractory}
  \to
  \mathrm{relative\ refractory}
  \to
  \mathrm{threshold\ reset}.
\]

Its period $\tau_{\mathrm{refrac}}$ determines the maximal firing rate.

\subsubsection{4. Channel-Gating-Zyklen$C_{\mathrm{gate}}$}

Ionenkanäle unterliegen zyklischen Zustandsübergängen:
\[
  S_i \in \{\mathrm{open}, \mathrm{closed}, \mathrm{inactive}\}.
\]

Der Gating-Zyklus wird angetrieben durch:
\[
  g_i(t) = g_{\mathrm{max}} m(t)^p h(t)^q,
\]
Wo$m,h$erfüllen ihre eigenen Wiederholungsgleichungen.

Der komplette Zyklus:
\[
  S_i^{(\mathrm{open})}
  \to 
  S_i^{(\mathrm{inactive})}
  \to
  S_i^{(\mathrm{closed})}
  \to
  S_i^{(\mathrm{open})}.
\]

Gating cycles stabilise $C_{\mathrm{spike}}$ and $C_{\mathrm{sub}}$.

\subsubsection{5. Kalzium- und Second-Messenger-Zyklen$C_{\mathrm{Ca}}$}

Dynamik der Calciumkonzentration:
\[
  \mathrm{Ca}^{2+}(t)
  \leftrightarrows
  \mathrm{buffered\; Ca},\;
  \mathrm{pumped\; Ca},
\]
bilden Zyklen, die Freisetzung, Sequestrierung und Pumpen umfassen.

Stabilität erfordert:
\[
  T_{\mathrm{Ca}} < \Theta_{\mathrm{noise-elec}}.
\]

Diese Zyklen modulieren sowohl die Erregbarkeit als auch die synaptische Veränderung.

\subsubsection{6. Synaptische Potenzierungs-/Depressionszyklen$C_{\mathrm{syn}}$}

Anhaltende Aktivität führt zu einer zyklischen Modulation des synaptischen Gewichts:
\[
  w_{ij}(t+\tau) = w_{ij}(t),
\]
angetrieben von:
\[
  J_{\mathrm{pre}},\; J_{\mathrm{post}},\; [\mathrm{Ca}^{2+}],
\]
und Schwellenbedingungen für Plastizität:
\[
  T_{\mathrm{plastic}} > \Theta_{\mathrm{plastic}}.
\]

Diese Zyklen ermöglichen Gedächtnis und dauerhafte Muster.

\subsubsection{7. Zyklen auf Netzwerkebene$C_{\mathrm{net}}$(Kontinuum$K_5'$)}

Using the representability theorem for neural networks (the corresponding derivation note),
Ein Netzwerk bildet ein erweitertes Kontinuum$K_5'$mit Zyklen:

\[
  C_{\mathrm{rhythm}},
  \qquad
  C_{\mathrm{sync}},
  \qquad
  C_{\mathrm{burst}}.
\]

Sie entstehen bei der kollektiven Aktivierung von Feldern
$p_i(t)$Erfüllen Sie die Wiederholung:
\[
  p_i(t+\tau) = p_i(t),
\]
und Schwellenwertbedingungen:
\[
  T_{\mathrm{sync}} < \Theta_{\mathrm{front}},\qquad
  T_{\mathrm{burst}} < \Theta_{\mathrm{spatial}}.
\]

Diese Zyklen sind die Grundlage organisierter Schwingungen und
Protokognition in$K_5'$.

\subsubsection{Metriken von Zyklen auf \texorpdfstring{$K_5$}{K_5}}

Using the scoped cycle metrics (the corresponding derivation note):

\paragraph{Length.}
\[
  L(C)=\oint d\Omega(K_5).
\]

\paragraph{Efficiency.}
\[
  C_{\mathrm{eff}}=\frac{\oint J\cdot dA}{L(C)},
\]
Wo$A$beinhaltet:
\[
  A_{\mathrm{exc}},\;
  A_{\mathrm{ion}},\;
  A_{\mathrm{syn}},\;
  A_{\mathrm{Ca}}.
\]

\paragraph{Stability.}
\[
S(C) =
    \min\limits_{t\in C}
      d\bigl(\Omega(K_5),\partial\Omega(K_5)\bigr).
\]

\paragraph{Weight.}
\[
w(C) =
    F\bigl(C_{\mathrm{eff}}, S(C), P_{\mathrm{ion}},
      P_{\mathrm{Ca}}, w_{ij}\bigr).
\]

\subsubsection{Zusammenbruch von \texorpdfstring{$C(K_5)$}{C(K_5)}}

According to the corresponding derivation note, collapse occurs if:

\[
  \Theta_{\mathrm{exc}},\;
  \Theta_{\mathrm{spike}},\;
  \Theta_{\mathrm{noise-elec}},\;
  \Theta_{\mathrm{front}},\;
  \Theta_{\mathrm{spatial}},\;
  \Theta_{\mathrm{refrac}}
\]
werden verletzt.

Dann:
\[
  C_j \to \varnothing,
\qquad
  \Omega(K_5)=\varnothing,
\qquad
  k_5\to 0.
\]

\subsubsection{Kontinuität von \texorpdfstring{$K_4$}{K_4} zu \texorpdfstring{$K_5$}{K_5}}

From Polish~W3 (the corresponding derivation note), the transition is $C^1$-continuous:
\[
  C_{\mathrm{exc}}^{(K_4)}
   \longrightarrow
  C_{\mathrm{spike}}^{(K_5)},
\]
mit sanftem Entstehen von:
\[
  A_{\mathrm{exc}},\quad
  \Theta_{\mathrm{spike}},\quad
  C_{\mathrm{refrac}}.
\]

Dieser sanfte Übergang gewährleistet die Lebensfähigkeit der ersten Nervenzellen.
